\subsection{Sparse covering scales and explicit irregular examples}
\label{sparsemoran:section}

The coherent comparison can use selected levels of one fixed positive
approximation sequence. This was already available in the block form of
the argument. The open family below and its closed endpoint were already
constructed in the earlier selected-scale research notes; the endpoint
there used a square-root gap. Here we give a complete proof and use the
weaker logarithmic gap at the endpoint. We record the estimate before
constructing the examples. The
resulting class goes beyond the preceding all-scale covering assumptions.
It does not assert a new sufficient inequality depending only on
Hausdorff and packing dimensions.

\begin{proposition}[A selected-scale covering criterion]
\label{sparsemoran:criterion}
Let \(\mu,\nu\) be compactly supported planar probabilities with
positively separated supports, where \(\mu\) is \(d\)-Frostman,
\(1<d<2\), and \(\nu\) has a Frostman exponent greater than one.
Let \(q>1\) be the averaged angular exponent and put
\(\eta_q=(q-1)/(2q-1)\). Suppose increasing integers \(n_j\to\infty\)
satisfy
\begin{equation}\label{sparsemoran:selectedcount}
 N_{n_j}\le C2^{dn_j},
\end{equation}
where \(N_n\) counts occupied source dyadic squares.
If, for some fixed \(0<\gamma<(d-1)/2\),
\begin{equation}\label{sparsemoran:series}
 \sum_j2^{-2\gamma\eta_q(2n_j-n_{j+1})}<\infty,
\end{equation}
then the raw joint distance law is absolutely continuous.
In particular it suffices that
\(n_{j+1}\le(2-\epsilon)n_j\) eventually for some \(\epsilon>0\).
It also suffices that \(n_{j+1}/n_j\to2\) and
\(2n_j-n_{j+1}\ge c\log_2(n_j+1)\) eventually, with \(c>0\).
\end{proposition}

\begin{proof}
Use the same chosen centers and positive affine densities \(f_n\)
as in Section~\ref{sec:coherenttree}. They do not depend on the
selected levels. Compare arbitrary \(n<m\), writing
\(r=2^{-n}\) and \(\delta=2^{-m}\). For a depth-\(m\) cell \(Q\)
inside its depth-\(n\) ancestor \(P\), both affine maps act on the
same conditional source \(\mu_Q\), recentered at \(x_Q\).
Their angles differ by \(O(r)\), and their scalar translation
constants differ by \(O(r^2)\).

After restricting both angular marginals to density at most \(A\),
use the angular estimate in Lemma~\ref{lem:coherentangular} and its
translation estimate with an arbitrary shift bound \(h>0\).
The latter follows from its same Plancherel proof, using
\(\lvert e^{i\xi v}-1\rvert\le C\lvert\xi h\rvert^\gamma\)
for \(\lvert v\rvert\le h\); it is not restricted to shifts smaller
than the square of the source diameter. With \(h=Cr^2\), these
estimates give an integrated squared
second-norm difference at most
\[
 CA\bigl(r^{2\gamma}\delta^{2\gamma}+r^{4\gamma}\bigr)
 I_{1+2\gamma}(\mu_Q)
 \le CA r^{4\gamma}I_{1+2\gamma}(\mu_Q).
\]
Each of the two conditional affine densities is supported in an
interval of length \(C\delta\). Their difference is supported in
the union of these intervals, whose total Lebesgue measure is
\(O(\delta)\), regardless of the gap between them. Thus the
first-norm Cauchy--Schwarz factor is \(\delta^{1/2}\);
a single containing interval is unnecessary.

Summing the child weights and the discarded pin masses exactly as
in Proposition~\ref{prop:coherentcomparison} gives
\begin{equation}\label{sparsemoran:block}
 \|f_m-f_n\|_1
 \le C(AZ_{n,m})^{1/2}+CBA^{1-q},\qquad
 Z_{n,m}=\delta r^{4\gamma}
       \sum_{Q\text{ at depth }m}p_QI_{1+2\gamma}(\mu_Q).
\end{equation}
The parent angular costs still sum correctly because all depth-\(m\)
descendants of a given depth-\(n\) parent have total mass equal
to that parent's mass. No restriction on \(m/n\) was used in
\eqref{sparsemoran:block}.

The conditional Frostman energy bound yields
\[
 I_{1+2\gamma}(\mu_Q)
 \le C\delta^{d-1-2\gamma}/p_Q.
\]
At \(m=n_{j+1}\), equation~\eqref{sparsemoran:selectedcount} therefore
gives
\[
 Z_{n_j,n_{j+1}}
 \le C r^{4\gamma}\delta^{d-2\gamma}N_{n_{j+1}}
 \le C2^{-2\gamma(2n_j-n_{j+1})}.
\]
Optimizing \(A\) in \eqref{sparsemoran:block} bounds the first-norm
difference by \(C B^{1/(2q-1)}Z_{n_j,n_{j+1}}^{\eta_q}\).
The assumed series makes \(f_{n_j}\) Cauchy in the joint first
norm. Its nonnegative mass-one limit is the actual joint law,
because the affine approximation error is uniformly
\(O(2^{-2n_j})\).

For the first stated sufficient condition the summands are bounded
by an exponential in \(-n_j\), and \(n_j\ge j\).
For the second, the levels eventually grow at least geometrically,
while the summands are bounded by a negative fixed power of \(n_j\).
Both series converge.
\end{proof}

\begin{theorem}[A class beyond the all-scale condition]
\label{sparsemoran:theorem}
For every
\begin{equation}\label{sparsemoran:range}
 1<d<D\le\frac{4d}{d+2}<2,
\end{equation}
there are a compact \(K\subset[0,1]^2\) and a \(d\)-Frostman
probability \(\mu\) on \(K\) such that
\[
 0<\mathcal H^d(K)<\infty,\qquad \HH K=d,\qquad \PP K=D,
\]
and every subset \(S\subset K\) obeys
\begin{equation}\label{sparsemoran:hereditary}
 \HH S\le\frac dD\,\PP S.
\end{equation}
The raw joint distance law with source \(\mu\) and any compactly
supported pin Frostman probability of exponent greater than one
is absolutely continuous. In particular
\[
 |\Delta_y(K)|>0\qquad\text{for \(\mu\)-almost every }y\in K.
\]
The endpoint \(D=4d/(d+2)\) is included.

If additionally \(d\le d_\alpha=(7-\sqrt7)/4\) and \(D\ge2d-1\),
no subset \(S\subset K\) with \(\HH S>1\) satisfies
\(\PP S<B_H(\HH S)\). No subset of \(K\) has positive equal
Hausdorff and packing dimensions.
\end{theorem}

\begin{proof}
We give the complete branching construction, allowing both the strict
interior and the endpoint of \eqref{sparsemoran:range}. Set
\[
 \ell(n)=
 \begin{cases}
  0,&D<4d/(d+2),\\
  A\log_2(n+1),&D=4d/(d+2),
 \end{cases}
 \qquad U(n)=Dn-\ell(n),
\]
where \(A>0\) is arbitrary in the endpoint case.
Choose a large \(a_0\) so that
\[
 U(n)-dn>10,\qquad 0<U(n+1)-U(n)<2
 \quad(n\ge a_0).
\]
This is possible since \(d<D<2\).
Start with \(F(a_0)=2\lceil da_0/2\rceil\), using an initial
integer path from \(F(0)=0\) with increments zero or two.

Starting from \(a_{j-1}\), increase \(F\) with slope two until
the first integer \(b_j\) with \(F(b_j)\ge U(b_j)\). Then keep
\(F\) constant until the first integer \(a_j>b_j\) with
\(da_j\ge F(b_j)\). These stopping times are finite since
\(D<2\). The first-crossing and first-catching properties give
\begin{equation}\label{sparsemoran:stops}
 \begin{gathered}
 U(b_j)\le F(b_j)<U(b_j)+2,\\
 F(a_j)=F(b_j),\qquad da_j-d<F(a_j)\le da_j.
 \end{gathered}
\end{equation}
Throughout the construction,
\begin{equation}\label{sparsemoran:barriers}
 dn-2\le F(n)\le Dn-\ell(n)+2
 \quad(n\ge a_0).
\end{equation}
Indeed the growing slope exceeds \(d\), the flat phase stops within
one step of the lower line, and \(U\) is increasing. Both phases
have positive length, so their endpoints tend to infinity.

At every dyadic depth use the same rule in every occupied square:
retain all four children when \(F\) increases by two, and only the
lower-left child when it stays constant. Equal conditional weights
give a compact \(K\) and a probability \(\mu\), with \(2^{F(n)}\)
occupied depth-\(n\) cells of mass \(2^{-F(n)}\).
There are infinitely many common zero digits. Thus an admissible
coordinate tail cannot be all ones, so half-open dyadic coding is
unambiguous; each fixed cylinder's descendants remain below its
upper edges.

The lower barrier and bounded overlap of balls with dyadic squares
give
\[
 \mu(B(x,2^{-n}))\le C2^{-F(n)}\le C2^{-dn}.
\]
Thus \(\mu\) is \(d\)-Frostman, and \(\mathcal H^d(K)>0\).
At the low endpoints,
\[
 2^{F(a_j)}(\sqrt2\,2^{-a_j})^d\le2^{d/2},
\]
which proves \(\mathcal H^d(K)<\infty\) and \(\HH K=d\).
The upper barrier gives \(\overline{\dim}_{\mathrm B}K\le D\).
At the peaks \(F(b_j)/b_j\to D\), and
\(\sup_x\mu(B(x,2^{-b_j}))\le C2^{-F(b_j)}\).
Every subset of positive \(\mu\) outer measure therefore has upper
box dimension at least \(D\). Applying the countable-cover
characterization of packing dimension proves \(\PP K=D\).

Since \(\ell(n)=o(n)\), the flat-phase identities give
\[
 \frac{a_j}{b_j}\longrightarrow\frac Dd.
\]
For any \(S\subset K\), cover \(S\) countably by subsets \(S_i\)
whose upper box dimensions are at most \(\PP S+\epsilon\).
At depth \(b_j\), each fixed \(S_i\) meets at most
\(C2^{(\PP S+2\epsilon)b_j}\) cells for large \(j\).
There is no branching between \(b_j\) and \(a_j\), so the same
number of depth-\(a_j\) cells covers \(S_i\).
The Hausdorff covering sums give
\[
 \HH S_i\le\frac dD(\PP S+2\epsilon).
\]
Countable stability and \(\epsilon\downarrow0\) prove
\eqref{sparsemoran:hereditary}, without any measurability assumption
on \(S\).

It remains to verify that the low endpoints are frequent enough
for Proposition~\ref{sparsemoran:criterion}. Write
\(a=a_{j-1}\), \(b=b_j\), and \(m=a_j\). The growing phase and
\eqref{sparsemoran:stops} imply
\[
 (2-D)b=(2-d)a-\ell(b)+O(1),\qquad
 dm=Db-\ell(b)+O(1).
\]
In particular \(b/a\to(2-d)/(2-D)\). Eliminating \(b\) gives
\begin{equation}\label{sparsemoran:recurrence}
 m=\kappa a-\frac{2}{d(2-D)}\ell(b)+O(1),
 \qquad
 \kappa=\frac{D(2-d)}{d(2-D)}.
\end{equation}
If \(D<4d/(d+2)\), then \(1<\kappa<2\), so the low endpoints
satisfy \(a_j\le(2-\epsilon)a_{j-1}\) eventually for some
\(\epsilon>0\).

At \(D=4d/(d+2)\), one has \(\kappa=2\). Since
\(\log_2(b+1)=\log_2(a+1)+O(1)\), the logarithmic choice gives
\begin{equation}\label{sparsemoran:endpointgap}
 2a-m=\frac{2A}{d(2-D)}\log_2(a+1)+O(1).
\end{equation}
Also \(m/a\to2\), so the low endpoints grow at least geometrically.
Thus they satisfy the second sufficient condition in
Proposition~\ref{sparsemoran:criterion}, for every \(q>1\) and every
fixed legal \(\gamma>0\). At all low endpoints the occupied count
is at most \(2^{da_j}\), by \eqref{sparsemoran:stops}.
The proposition proves raw joint absolute continuity for separated
supports.

To remove separation, cover the complement of the source--pin
diagonal by countably many products of separated compact balls.
Positive normalized restrictions preserve the source Frostman
bound and the selected-scale occupied counts, and preserve the
pin Frostman property. Each restricted joint law is absolutely
continuous. The diagonal has zero pair mass because \(\mu\) is
nonatomic. The countable-cover argument therefore proves absolute
continuity of the full joint law, including self-pins.

Finally, suppose \(d\le d_\alpha\), \(D\ge2d-1\), and
\(h=\HH S>1\). Then \(h\le d\), \(B_H(h)=2h-1\), and
\[
 \PP S\ge(D/d)h,\qquad
 (D/d)h-(2h-1)\ge1-h/d\ge0.
\]
This excludes the previous strict criterion on every such subset.
Equation~\eqref{sparsemoran:hereditary}, with \(d/D<1\), also
excludes every positive-dimensional regular subset.
\end{proof}

For a concrete instance choose
\begin{equation}\label{sparsemoran:example}
 d=\frac{21}{20},\qquad D=\frac65.
\end{equation}
Here \(D<4d/(d+2)=84/61\), the low-endpoint ratio is
\(\kappa=19/14<2\), and \(B_H(d)=11/10<D\).
Thus this fixed set has positive-length self-pinned distances for
almost every natural pin, and every subset of Hausdorff dimension
greater than one fails the preceding strict dimension criterion.
The constructed class also reaches \(D=84/61\) at the same
Hausdorff dimension by \eqref{sparsemoran:endpointgap}.

The favorable information is the recurrence of the low covering
scales, not the pair of dimension numbers alone. In the zero-gauge
interior construction \(F(b_j)=Db_j+O(1)\) at infinitely many peaks.
For \(D\ge2d-1\), the all-scale summands based on
\(2^{-n(2d-1)}N_n\) consequently fail to tend to zero there.
The sparse proof demonstrates why failure of that numerical
all-scale criterion is not an obstruction to the actual distance law.
